% ============================================================================
%  Supplementary Material
%  Regime-Adaptive Pairs Trading with Robust Kalman Filtering
% ============================================================================
\documentclass[conference]{IEEEtran}

\usepackage[utf8]{inputenc}
\usepackage[T1]{fontenc}
\usepackage{amsmath,amssymb,amsfonts}
\usepackage{graphicx}
\usepackage{booktabs}
\usepackage{multirow}
\usepackage{hyperref}
\usepackage{xcolor}
\usepackage{tabularx}
\usepackage{longtable}
\usepackage{float}

\hypersetup{
    colorlinks=true,
    linkcolor=blue!70!black,
    citecolor=blue!70!black,
    urlcolor=blue!70!black,
}

\begin{document}

\title{Supplementary Material:\\
Regime-Adaptive Pairs Trading with Robust Kalman Filtering}

\author{
\IEEEauthorblockN{Anonymous Authors}
\IEEEauthorblockA{University}
}

\maketitle

\noindent This supplement provides (A)~complete hyperparameter settings,
(B)~extended sensitivity analysis, (C)~per-pair ablation tables,
(D)~walk-forward fold-by-fold results, (E)~complete feature importance
rankings, (F)~computational cost analysis, and (G)~pair discovery pipeline
details.

% ====================================================================
\section{A. Complete Hyperparameter Settings}
\label{sec:hyperparams}

All parameters are centralised in \texttt{Research/config\_experiments.yaml}
and reproduced here for convenience. No parameter was tuned on
out-of-sample data.

\subsection{A.1 Strategy Parameters}

\begin{table}[h]
\centering
\caption{Strategy hyperparameters by regime.}
\label{tab:supp_strategy}
\small
\begin{tabular}{lrrr}
\toprule
Parameter & Normal & Volatile & Crisis \\
\midrule
Entry z-score & 1.5 & 2.0 & $\infty$ (no entry) \\
Exit z-score & 0.5 & 0.3 & $\infty$ (exit all) \\
Min hold days & 3 & 3 & 0 \\
Max hold days & 30 & 30 & 0 \\
Stop-loss multiplier & 2.0$\times$ & 2.0$\times$ & --- \\
Z-score window & 40 & 40 & --- \\
Position sizing range & $[0.3, 1.0]$ & $[0.3, 1.0]$ & 0 \\
\bottomrule
\end{tabular}
\end{table}

\subsection{A.2 Kalman Filter Parameters}

\begin{table}[h]
\centering
\caption{Kalman filter hyperparameters.}
\label{tab:supp_kalman}
\small
\begin{tabular}{lr}
\toprule
Parameter & Value \\
\midrule
Base process noise $q_0$ ($\delta$) & $1 \times 10^{-4}$ \\
Observation noise $R$ (\texttt{var\_e}) & $1 \times 10^{-3}$ \\
Process noise offset (\texttt{var\_eta}) & $1 \times 10^{-4}$ \\
Crisis noise multiplier & 5$\times$ \\
Volatile noise multiplier & 2$\times$ \\
MAD outlier threshold & 3$\times$ \\
Innovation history size & 50 \\
Initial hedge ratio & $\beta_0 = Y_0 / X_0$ \\
Initial state covariance & $C_0 = 1.0$ \\
\bottomrule
\end{tabular}
\end{table}

Note: Eq.~(4) in the paper and the implementation are consistent for
Kalman process noise: multipliers $(1,2,5)$ for (Normal, Volatile, Crisis).
Signal-generation uses separate regime-dependent entry/exit thresholds
(Eq.~(10): $(1.5,0.5)$ Normal, $(2.0,0.3)$ Volatile, and no new entries in Crisis).

\subsection{A.3 Random Forest Regime Classifier}

\begin{table}[h]
\centering
\caption{RF classifier hyperparameters.}
\label{tab:supp_rf}
\small
\begin{tabular}{lr}
\toprule
Parameter & Value \\
\midrule
Number of estimators & 300 \\
Maximum depth & 6 \\
Minimum samples split & 40 \\
Minimum samples leaf & 20 \\
Random state & 42 \\
\midrule
\multicolumn{2}{l}{\textit{Regime label thresholds (vol quantiles)}} \\
Crisis threshold ($\sigma_{20}$) & 90th percentile \\
Volatile threshold ($\sigma_{20}$) & 70th percentile \\
\midrule
\multicolumn{2}{l}{\textit{Alternative thresholds (drawdown-based)}} \\
Crisis drawdown & $< -20\%$ \\
Volatile drawdown & $< -10\%$ \\
Crisis volatility & $> 0.35$ \\
Volatile volatility & $> 0.22$ \\
\bottomrule
\end{tabular}
\end{table}

\subsection{A.4 Features for Regime Classifier (21 total)}

\begin{table}[h]
\centering
\caption{Complete feature list for the RF regime classifier.}
\label{tab:supp_features}
\small
\begin{tabular}{rll}
\toprule
\# & Feature Name & Description \\
\midrule
1 & $\sigma_5$ & 5-day rolling volatility \\
2 & $\sigma_{10}$ & 10-day rolling volatility \\
3 & $\sigma_{20}$ & 20-day rolling volatility \\
4 & $\sigma_{40}$ & 40-day rolling volatility \\
5 & $\sigma_{60}$ & 60-day rolling volatility \\
6 & Vol trend & $\sigma_5 / \sigma_{60}$ ratio \\
7 & Vol acceleration & $\Delta(\sigma_5 / \sigma_{60})$ \\
8 & DD 5-day & 5-day trailing drawdown \\
9 & DD 10-day & 10-day trailing drawdown \\
10 & DD 20-day & 20-day trailing drawdown \\
11 & DD 40-day & 40-day trailing drawdown \\
12 & Return 5-day & 5-day cumulative return \\
13 & Return 10-day & 10-day cumulative return \\
14 & Return 20-day & 20-day cumulative return \\
15 & Trend & 40-day SMA of price returns \\
16 & Trend strength & $|$Trend$|$ \\
17 & Range expansion & Ratio of recent to historical range \\
18 & Pair correlation & 40-day rolling correlation \\
19 & Pair corr change & $\Delta$(pair correlation) \\
20 & Pair volatility & $\sigma_{20}$ of the spread \\
21 & Pair spike & Binary: spread $> 2\sigma$ from mean \\
\bottomrule
\end{tabular}
\end{table}

\subsection{A.5 Transaction Cost Model}

\begin{table}[h]
\centering
\caption{Transaction cost components.}
\label{tab:supp_costs}
\small
\begin{tabular}{lrl}
\toprule
Component & Value & Notes \\
\midrule
Slippage & 5 bps & Per side, proportional to notional \\
Commission & \$1.00 & Per trade (round-trip: \$2) \\
Market impact & $0.1 \sigma_t \sqrt{V/\text{ADV}}$ & Almgren-Chriss style \\
\bottomrule
\end{tabular}
\end{table}

All costs are deducted \emph{inside} the backtest loop at each trade
execution, not applied post-hoc.

\subsection{A.6 Evaluation Parameters}

\begin{table}[h]
\centering
\caption{Statistical evaluation hyperparameters.}
\label{tab:supp_eval}
\small
\begin{tabular}{lr}
\toprule
Parameter & Value \\
\midrule
Walk-forward min training & 504 days ($\approx$ 2 years) \\
Walk-forward step size & 63 days ($\approx$ 1 quarter) \\
Walk-forward test window & 126 days ($\approx$ 6 months) \\
Walk-forward min folds & 8 \\
Walk-forward total folds & 32 \\
Monte Carlo simulations & 10,000 \\
Bootstrap resamples & 10,000 \\
Bootstrap block size & 5 \\
Confidence level & 95\% \\
FDR $\alpha$ (Benjamini-Hochberg) & 0.05 \\
\bottomrule
\end{tabular}
\end{table}

% ====================================================================
\section{B. Extended Sensitivity Analysis}
\label{sec:sensitivity}

\subsection{B.1 Entry/Exit Z-Score Grid}

The full 7$\times$5 parameter grid tested (35 configurations per pair):

\begin{table}[h]
\centering
\caption{Sensitivity grid: entry z $\times$ exit z values tested.}
\label{tab:supp_grid}
\small
\begin{tabular}{ll}
\toprule
Parameter & Values \\
\midrule
Entry z-score & 1.0, 1.25, 1.5, 1.75, 2.0, 2.5, 3.0 \\
Exit z-score & 0.2, 0.3, 0.5, 0.7, 1.0 \\
Z-score window & 30, 40, 60 \\
\bottomrule
\end{tabular}
\end{table}

\subsection{B.2 BAC/PNC Sensitivity Results (Selected)}

\begin{table}[h]
\centering
\caption{BAC/PNC OOS Sharpe across selected entry/exit combinations.}
\label{tab:supp_bac_sens}
\small
\begin{tabular}{l|rrrrr}
\toprule
Entry $\backslash$ Exit & 0.2 & 0.3 & 0.5 & 0.7 & 1.0 \\
\midrule
1.0 & $-$0.31 & $-$0.28 & $-$0.22 & $-$0.19 & $-$0.15 \\
1.5 & $-$0.18 & $-$0.13 & $-$0.13 & $-$0.09 & +0.02 \\
2.0 & $-$0.06 & $-$0.04 & +0.08 & +0.15 & +0.22 \\
2.5 & +0.03 & +0.07 & +0.18 & +0.24 & +0.31 \\
3.0 & +0.11 & +0.15 & +0.22 & +0.28 & +0.19 \\
\bottomrule
\end{tabular}
\end{table}

Sharpe $> 0$ is concentrated in the upper-right corner (high entry, high exit),
where the strategy executes very few trades. The profitable region covers
approximately 22\% of the grid --- below the 60\% threshold we consider robust.

% ====================================================================
\section{C. Per-Pair Ablation Details}
\label{sec:ablation_detail}

\subsection{C.1 Ablation Configurations}

\begin{table}[h]
\centering
\caption{Nine ablation configurations and what each removes.}
\label{tab:supp_configs}
\small
\begin{tabular}{rl}
\toprule
\# & Configuration \\
\midrule
1 & Full system (all components active) \\
2 & Remove regime classification ($\rightarrow$ always Normal) \\
3 & Remove adaptive Kalman ($\rightarrow$ fixed OLS hedge) \\
4 & Remove MAD outlier detection \\
5 & Remove dynamic z-score window \\
6 & Remove stop-loss mechanism \\
7 & Remove minimum holding period \\
8 & Remove transaction costs (post-hoc only) \\
9 & Remove volatile-regime parameters ($\rightarrow$ Normal thresholds) \\
\bottomrule
\end{tabular}
\end{table}

\subsection{C.2 WFC/MS Ablation Results}

\begin{table}[h]
\centering
\caption{WFC/MS ablation results (OOS period).}
\label{tab:supp_wfc_abl}
\small
\begin{tabular}{lrrr}
\toprule
Config & OOS Sharpe & OOS Ret.~\% & Max DD~\% \\
\midrule
Full system & $-$0.045 & $-$2.71 & $-$37.5 \\
No regime & $-$0.028 & $-$0.99 & $-$35.2 \\
No Kalman (OLS) & +0.097 & +5.23 & $-$28.1 \\
No MAD & $-$0.012 & $-$0.84 & $-$36.8 \\
No dynamic z-window & $-$0.051 & $-$3.12 & $-$38.1 \\
No stop-loss & $-$0.039 & $-$2.45 & $-$42.3 \\
No min hold & $-$0.058 & $-$3.51 & $-$39.7 \\
No costs & +0.031 & +1.87 & $-$36.4 \\
No volatile params & $-$0.042 & $-$2.55 & $-$37.0 \\
\bottomrule
\end{tabular}
\end{table}

\subsection{C.3 CVX/OXY Ablation Results}

\begin{table}[h]
\centering
\caption{CVX/OXY ablation results (OOS period).}
\label{tab:supp_cvx_abl}
\small
\begin{tabular}{lrrr}
\toprule
Config & OOS Sharpe & OOS Ret.~\% & Max DD~\% \\
\midrule
Full system & $-$0.613 & $-$90.59 & $-$95.4 \\
No regime & $-$0.571 & $-$46.48 & $-$92.1 \\
No Kalman (OLS) & $-$0.312 & $-$28.34 & $-$78.5 \\
No MAD & $-$0.598 & $-$88.23 & $-$94.8 \\
No dynamic z-window & $-$0.625 & $-$91.42 & $-$95.7 \\
No stop-loss & $-$0.641 & $-$92.87 & $-$96.2 \\
No min hold & $-$0.629 & $-$91.98 & $-$95.9 \\
No costs & $-$0.507 & $-$72.31 & $-$93.1 \\
No volatile params & $-$0.608 & $-$89.81 & $-$95.1 \\
\bottomrule
\end{tabular}
\end{table}

CVX/OXY shows uniformly negative results across all configurations,
confirming that the energy sector structural break overwhelms all
algorithmic improvements.

% ====================================================================
\section{D. Walk-Forward Fold-by-Fold Results}
\label{sec:walkforward_detail}

\subsection{D.1 BAC/PNC Walk-Forward (Representative Folds)}

\begin{table}[h]
\centering
\caption{BAC/PNC walk-forward: selected folds (32 total).}
\label{tab:supp_wf_bac}
\small
\begin{tabular}{rrlrr}
\toprule
Fold & Train Days & Test Period & Sharpe & Ret.~\% \\
\midrule
1 & 504 & 2017Q1 & +0.42 & +3.1 \\
4 & 693 & 2017Q4 & +0.28 & +1.8 \\
8 & 1008 & 2018Q4 & $-$0.15 & $-$1.2 \\
12 & 1260 & 2019Q4 & +0.51 & +4.2 \\
16 & 1511 & 2021Q1 & $-$0.08 & $-$0.6 \\
20 & 1762 & 2022Q1 & $-$0.34 & $-$2.8 \\
24 & 2013 & 2023Q1 & +0.12 & +0.9 \\
28 & 2264 & 2024Q1 & $-$0.22 & $-$1.7 \\
32 & 2516 & 2025Q1 & $-$0.11 & $-$0.8 \\
\midrule
\multicolumn{3}{l}{Average (all 32 folds)} & +0.083 & $-$0.78 \\
\multicolumn{3}{l}{Consistency (folds with Sharpe $> 0$)} & \multicolumn{2}{r}{50\%} \\
\bottomrule
\end{tabular}
\end{table}

\subsection{D.2 Walk-Forward Stability}

\begin{table}[h]
\centering
\caption{Walk-forward stability metrics across pairs.}
\label{tab:supp_wf_stability}
\small
\begin{tabular}{lrrrr}
\toprule
Pair & Sharpe $\bar{x}$ & Sharpe $\sigma$ & Slope & $p_{\text{slope}}$ \\
\midrule
BAC/PNC & +0.083 & 0.412 & $-$0.012 & 0.192 \\
WFC/MS & +0.096 & 0.387 & +0.005 & 0.614 \\
CVX/OXY & $-$0.275 & 0.562 & $-$0.018 & 0.087 \\
\bottomrule
\end{tabular}
\end{table}

No pair shows a statistically significant time trend in fold-level Sharpe
($p > 0.05$ for all), indicating the underperformance is persistent rather
than worsening.

% ====================================================================
\section{E. Complete Feature Importance Rankings}
\label{sec:features_detail}

\begin{table}[h]
\centering
\caption{Full 21-feature ranking by consensus importance (BAC/PNC).}
\label{tab:supp_feature_full}
\small
\begin{tabular}{rlrrr}
\toprule
Rank & Feature & SHAP & Gini & Permutation \\
\midrule
1 & $\sigma_{20}$ & 0.081 & 0.188 & 0.042 \\
2 & Pair volatility & 0.058 & 0.092 & 0.031 \\
3 & Pair spike & 0.045 & 0.078 & 0.028 \\
4 & $\sigma_{40}$ & 0.039 & 0.071 & 0.024 \\
5 & DD 20-day & 0.034 & 0.065 & 0.021 \\
6 & Vol trend & 0.031 & 0.058 & 0.019 \\
7 & $\sigma_{10}$ & 0.028 & 0.052 & 0.017 \\
8 & $\sigma_{60}$ & 0.025 & 0.048 & 0.015 \\
9 & Return 20-day & 0.022 & 0.041 & 0.013 \\
10 & DD 10-day & 0.019 & 0.037 & 0.011 \\
\midrule
11 & $\sigma_5$ & 0.017 & 0.034 & 0.010 \\
12 & Pair correlation & 0.015 & 0.031 & 0.009 \\
13 & Vol acceleration & 0.013 & 0.028 & 0.008 \\
14 & Return 10-day & 0.011 & 0.025 & 0.007 \\
15 & DD 40-day & 0.010 & 0.022 & 0.006 \\
16 & DD 5-day & 0.008 & 0.019 & 0.005 \\
17 & Return 5-day & 0.007 & 0.016 & 0.004 \\
18 & Trend & 0.006 & 0.014 & 0.003 \\
19 & Trend strength & 0.005 & 0.012 & 0.003 \\
20 & Range expansion & 0.004 & 0.010 & 0.002 \\
21 & Pair corr change & 0.003 & 0.008 & 0.002 \\
\bottomrule
\end{tabular}
\end{table}

The top-5 features account for 63\% of total SHAP importance. Pair-specific
features (ranks 2--3) are the most significant non-volatility features,
suggesting cross-asset information contributes to regime detection beyond
what market volatility alone provides.

% ====================================================================
\section{F. Computational Cost Analysis}
\label{sec:compute}

\begin{table}[h]
\centering
\caption{Runtime breakdown for full reproduction (MacBook Pro M2, 16 GB RAM).}
\label{tab:supp_runtime}
\small
\begin{tabular}{lrr}
\toprule
Component & Time (min) & \% of Total \\
\midrule
Data download (yfinance) & 2.5 & 7\% \\
Baseline comparison ($3 \times 7$) & 5.2 & 14\% \\
Ablation study ($3 \times 9$) & 6.8 & 19\% \\
Feature analysis (SHAP) & 3.1 & 9\% \\
Statistical tests (bootstrap) & 8.4 & 23\% \\
Walk-forward ($3 \times 32$ folds) & 4.7 & 13\% \\
Cointegration analysis & 1.2 & 3\% \\
Regime evaluation & 2.3 & 6\% \\
Attribution & 1.0 & 3\% \\
Figure generation & 0.8 & 2\% \\
Table generation & 0.2 & 1\% \\
\midrule
\textbf{Total} & \textbf{36.2} & \textbf{100\%} \\
\bottomrule
\end{tabular}
\end{table}

The \texttt{--quick} flag reduces bootstrap/Monte Carlo from 10,000 to 1,000
iterations, cutting the statistical tests phase from 8.4 to 0.8 minutes
(total: $\sim$28 min $\rightarrow$ $\sim$8 min).

\subsection{F.1 Memory Usage}

\begin{table}[h]
\centering
\caption{Peak memory usage by experiment.}
\label{tab:supp_memory}
\small
\begin{tabular}{lr}
\toprule
Component & Peak RAM (MB) \\
\midrule
Data loading (3 pairs + SPY) & 45 \\
RF training (300 trees) & 120 \\
SHAP computation & 310 \\
Bootstrap (10,000 resamples) & 180 \\
Walk-forward (32 folds) & 250 \\
Figure rendering & 150 \\
\midrule
Overall peak & $\sim$400 \\
\bottomrule
\end{tabular}
\end{table}

All experiments fit comfortably in 4 GB RAM. The Docker container
requires approximately 2.1 GB disk (Python 3.12 slim + dependencies).

% ====================================================================
\section{G. Pair Discovery Pipeline}
\label{sec:pairs}

The \texttt{Pair\_Discovery/auto\_find\_pairs.py} module screens for
candidate pairs using the following pipeline:

\begin{enumerate}
    \item Download daily prices for all stocks in the target sector(s)
    from Yahoo Finance (minimum 504 trading days).
    \item Compute pairwise Engle-Granger cointegration $p$-values on each
    pair's training-period data.
    \item Filter pairs with $p < 0.05$.
    \item Compute half-life of mean reversion via OLS on the spread
    ($\tau_{1/2} = \ln 2 / \hat{\theta}$); require $5 < \tau_{1/2} < 60$ days.
    \item Compute Hurst exponent; require $H < 0.5$.
    \item Rank remaining pairs by a composite score combining $p$-value,
    Hurst exponent, and correlation stability.
    \item Select top pairs (max 10 per sector).
\end{enumerate}

\begin{table}[h]
\centering
\caption{Discovered pairs and their in-sample quality metrics.}
\label{tab:supp_pairs}
\small
\begin{tabular}{llrrrr}
\toprule
Pair & Sector & ADF $p$ & Hurst $H$ & Half-life & Score \\
\midrule
BAC/PNC & Banking & 0.012 & 0.42 & 18.3 & 72 \\
WFC/MS & Financial & 0.031 & 0.44 & 22.1 & 65 \\
CVX/OXY & Energy & 0.008 & 0.38 & 14.7 & 78 \\
\bottomrule
\end{tabular}
\end{table}

All three pairs showed good in-sample cointegration properties, but
OOS performance degraded significantly (see main paper Table~1).
This discrepancy between in-sample quality and OOS performance is the
central challenge in pairs trading research.

% ====================================================================
\section{H. Software and Reproducibility}
\label{sec:software}

\begin{table}[h]
\centering
\caption{Software environment.}
\label{tab:supp_software}
\small
\begin{tabular}{lr}
\toprule
Component & Version \\
\midrule
Python & 3.12.9 \\
NumPy & 2.3.5 \\
pandas & 2.3.3 \\
SciPy & 1.16.3 \\
scikit-learn & 1.7.2 \\
statsmodels & 0.14.5 \\
SHAP & 0.51.0 \\
yfinance & 1.1.0 \\
matplotlib & 3.10.7 \\
seaborn & 0.13.2 \\
\bottomrule
\end{tabular}
\end{table}

\textbf{Reproducibility protocol:}
\begin{itemize}
    \item All random seeds fixed to 42 at top of every module.
    \item Pinned dependency versions in \texttt{requirements.txt}.
    \item Dockerfile provided for containerised reproduction.
    \item Single-command reproduction: \texttt{python Research/reproduce.py}.
    \item Intermediate results saved as CSV/JSON in \texttt{Research/results/}.
    \item Reproduction manifest (JSON) with timestamps, versions, and checksums.
\end{itemize}

\textbf{Limitations on exact reproducibility:}
\begin{itemize}
    \item Yahoo Finance data may be retrospectively adjusted (splits,
    dividends, corrections). The qualitative conclusions are robust to
    small data shifts, but exact numerical values may differ by
    $< 0.01$ Sharpe.
    \item Floating-point arithmetic is platform-dependent. Docker + pinned
    versions minimises but does not eliminate this.
\end{itemize}

\end{document}
